\section{Introduction à la statistique descriptive}

\begin{cadreobjectif}
La statistique descriptive permet :
\begin{itemize}
    \item d’organiser des données ;
    \item de résumer une série statistique ;
    \item de représenter graphiquement des informations ;
    \item d’interpréter des résultats dans un contexte réel.
\end{itemize}
\end{cadreobjectif}

\begin{cadretheorie}
La statistique descriptive est une branche des mathématiques qui consiste à :
\begin{itemize}
    \item collecter ;
    \item organiser ;
    \item représenter ;
    \item analyser ;
    \item interpréter
\end{itemize}
des données issues d’une situation réelle.
\end{cadretheorie}

\section{Vocabulaire statistique}

\begin{cadretheorie}
\textbf{Population :} ensemble des individus étudiés.

\medskip

\textbf{Individu :} élément appartenant à la population.

\medskip

\textbf{Échantillon :} partie de la population utilisée pour l’étude statistique.

\medskip

\textbf{Caractère statistique :} propriété étudiée sur les individus.
\end{cadretheorie}

\begin{cadreexemple}
Dans une école :

\begin{itemize}
    \item Population : tous les élèves ;
    \item Individu : un élève ;
    \item Caractère : la taille, l’âge ou le nombre d’heures d’étude.
\end{itemize}
\end{cadreexemple}

\section{Types de caractères statistiques}

\begin{cadretheorie}
On distingue deux grandes catégories de caractères statistiques :
\end{cadretheorie}

\begin{tableaularge}{|p{4cm}|p{5.5cm}|p{5.5cm}|}
\hline
\textbf{Type} & \textbf{Description} & \textbf{Exemple} \\
\hline

Qualitatif & Donnée non numérique & couleur des yeux \\
\hline

Quantitatif discret & Valeurs numériques isolées & nombre d’enfants \\
\hline

Quantitatif continu & Valeurs dans un intervalle & taille, masse \\
\hline

\end{tableaularge}

\section{Organisation des données}

\begin{cadretheorie}
Une série statistique peut être organisée dans un tableau contenant :
\begin{itemize}
    \item les valeurs observées ;
    \item les effectifs ;
    \item les fréquences ;
    \item les fréquences cumulées.
\end{itemize}
\end{cadretheorie}

\begin{cadreexemple}
Exemple : nombre de livres lus pendant un mois.
\end{cadreexemple}

\begin{center}
\begin{tabular}{|c|c|}
\hline
Nombre de livres & Effectif \\
\hline
0 & 3 \\
1 & 5 \\
2 & 7 \\
3 & 4 \\
\hline
\end{tabular}
\end{center}

\section{Importance des statistiques}

\begin{cadresynthese}
Les statistiques sont utilisées dans :
\begin{itemize}
    \item les sciences ;
    \item l’économie ;
    \item la médecine ;
    \item le sport ;
    \item les enquêtes d’opinion ;
    \item l’enseignement.
\end{itemize}

Une bonne analyse statistique nécessite toujours :
\begin{itemize}
    \item une lecture critique ;
    \item une interprétation dans le contexte ;
    \item une représentation adaptée.
\end{itemize}
\end{cadresynthese}